% \iffalse meta-comment
%
% hit-caption.dtx 是各个类共用的双语题注
%
% \fi
%
% \section{共享双语题注}
%
% \changes{v3.2a}{2026/09/14}{双语题注抽出来共用：英文题注是 \cs{caption} 的尾随可选参数，
%   \cs{bicaption} 降为兼容写法；短标题接到 \pkg{ccaption} 的第二个参数（原先被当成
%   \cs{label}），题注后补 \cs{@currentlabel}；\option{styles/bilingual-caption}
%   默认按学位只有博士开（规范“硕士学位论文只用中文，博士学位论文用中、英两种文字”）；
%   浮动体的四个间距统一交给 \cs{captionsetup}，\option{style}/\option{float-sep} 可改；
%   加 \cs{figurenote}；题注的字号行距从 \option{styles} 表取；\cs{caption*} 排不编号的题注，
%   也认尾随的英文题}
%
% 中英文对照的题注不换命令，用户不必记住“这张图要不要双语”再挑命令：
% \begin{latex}
% \caption{中文题}
% \caption{中文题}[English caption]
% \caption[目录用短标题]{中文题}[English caption]
% \end{latex}
% 给了尾随的 \oarg{英文题} 就排双语，不给就单语。目录用的短标题仍是 \cs{caption}
% 原本那个前置可选参数，两者位置不同，不会混。
%
% \option{styles/bilingual-caption} 可以统一关闭双语题注，同一份稿子投不同用途时
% 不必逐条删；默认按学位，见下面。
%
% \cs{bicaption} 仍然可用：\pkg{ccaption} 的五参数原形与本模板之前的三、四参数写法
% 都接得住，只是各提示一次，v4 移除。
%
%    \begin{macrocode}
%<*shared-caption>
%    \end{macrocode}
%
% 浮动体的间距照 \pkg{thuthesis} 的办法收进类里。三件事：
%
% 一、四个位置同一个值。浮动体落在正文中间、页顶页底、两个浮动体之间、整页浮动，
% \hologo{LaTeX} 内核给的是 8.5pt／8.5pt／12pt／10pt 四个样。
%
% 二、题注与图表之间的空白由 \cs{captionsetup} 的 \texttt{skip} 给，配上
% \texttt{tableposition} 与 \texttt{figureposition}，\pkg{caption} 自己知道表题在上
% 该把空白放下面、图题在下该放上面。旧版面取 6bp，与 v3.1e 示例里手写的
% \cs{vspace}\{0.5em\} 等价；新版面在 \file{src/flow/hit-final-float.dtx} 里归零，
% 范例里表题到表顶线只有表题自己那一行的行深。\pkg{caption} 由 \pkg{subcaption}
% 带进来，一直都在。
%
% 三、\hologo{LaTeX} 内核里浮动体与正文的前后距离本身不对称，光把长度设成一样没用，要打
% 三个补丁。出处见 \pkg{thuthesis} 的 \texttt{tuna/thuthesis\#614}。这段必须退出
% \pkg{expl3} 写：\cs{patchcmd} 把宏体按当前 catcode 重扫，\texttt{\string\vskip\space
% \string\intextsep} 里那个空格在 \pkg{expl3} 下会被吞，补丁就对不上（CONTRIBUTING §4.7）。
%
% 放在 \cs{AtBeginDocument} 里。\cs{AtEndOfClass} 那个时机太早：\cs{@addtocurcol}
% 那时还不是最终那一版，三个补丁全都找不到目标（实测报“没打上”），到
% \cs{begin\{document\}} 时 \cs{ifpatchable} 才为真。长度也一样，晚设才不会被覆盖。
%    \begin{macrocode}
\ExplSyntaxOff
\AtBeginDocument{%
  \setlength{\floatsep}{12bp \@plus 2bp \@minus 2bp}%
  \setlength{\textfloatsep}{12bp \@plus 2bp \@minus 2bp}%
  \setlength{\intextsep}{12bp \@plus 2bp \@minus 2bp}%
  \setlength{\@fptop}{0bp \@plus 1fil}%
  \setlength{\@fpsep}{12bp \@plus 2fil}%
  \setlength{\@fpbot}{0bp \@plus 1fil}%
  \captionsetup{skip=6bp,tableposition=top,figureposition=bottom}%
  \patchcmd{\@addtocurcol}%
    {\vskip \intextsep}%
    {\edef\hit@save@first@penalty{\the\lastpenalty}\unpenalty
     \ifnum \lastpenalty = \@M
       \unpenalty
     \else
       \penalty \hit@save@first@penalty \relax
     \fi
     \ifnum\outputpenalty <-\@Mii
       \addvspace\intextsep
       \vskip\parskip
     \else
       \addvspace\intextsep
     \fi}%
    {}{\ClassWarningNoLine{hithesis}{补丁 \string\@addtocurcol 之一没打上}}%
  \patchcmd{\@addtocurcol}%
    {\vskip\intextsep \ifnum\outputpenalty <-\@Mii \vskip -\parskip\fi}%
    {\ifnum\outputpenalty <-\@Mii
       \aftergroup\vskip\aftergroup\intextsep
       \aftergroup\nointerlineskip
     \else
       \vskip\intextsep
     \fi}%
    {}{\ClassWarningNoLine{hithesis}{补丁 \string\@addtocurcol 之二没打上}}%
  \patchcmd{\@getpen}{\@M}{\@Mi}%
    {}{\ClassWarningNoLine{hithesis}{补丁 \string\@getpen 没打上}}%
}
\ExplSyntaxOn
%    \end{macrocode}
%
% \begin{macro}{\hit@setup@caption}
% 接管 \cs{caption} 要等到 \texttt{begindocument/end}。\pkg{caption}（由
% \pkg{subcaption} 带进来）和 \pkg{ccaption} 都在 \cs{AtBeginDocument} 里重定义
% \cs{caption}，在类文件里接管会被它们覆盖掉——踩过一次，接管看似成功，
% \cs{meaning}\cs{caption} 到了正文里已经是 \pkg{caption} 的定义了。
% \texttt{begindocument/end} 排在所有 \cs{AtBeginDocument} 之后。
%
% \cs{\_\_hit\_bicaption:w} 存的是 \pkg{ccaption} 的原形，这个在类文件里存就行；
% \cs{hit@original@caption} 存的是接管前那条完整的 \cs{caption} 链。
%    \begin{macrocode}
\bool_new:N \g__hit_bilingual_caption_bool
\bool_new:N \g__hit_bilingual_caption_auto_bool
\bool_gset_true:N \g__hit_bilingual_caption_auto_bool
%    \end{macrocode}
%
% 双语题注是题注自己的属性，归 \option{styles}。写成族级键而不是挂在
% \option{caption} 底下，是因为 \option{caption} 是个 meta，会把整包取值转给
% \option{caption-figure} 与 \option{caption-table} 两个目标，而这是一个开关、
% 不分目标。
%
% 默认按学位：博士与博士后开，本科硕士关——硕士写了尾随的英文题也不排，规范
% 说硕士只用中文。博士后没有自己的图表条款，出站报告照博士的规矩排，v3.1e
% 以来也一直是双语。报告没有博士后这一档，那边只查博士，标志位不存在就当假。
% 全员开（“写了就出、不写就不出”）等于把合规交给用户自己删英文题；按学位给
% 默认之后合规是默认态，确要双语的硕士写 \option{bilingual-caption}\texttt{=true}
% 显式设为 \texttt{true} 即可覆盖默认值。三态里的 \texttt{auto} 就是按学位查表，
% 与 \option{openright} 同一个路数。
%
% 这个键挂在本类的 \texttt{hit / styles} 族上，与样式写在同一个 \option{styles} 里，
% 见 \file{hit-docxlike.dtx}。
%    \begin{macrocode}
\keys_define:nn { hit / styles }
  {
    bilingual-caption .choice: ,
    bilingual-caption / true .code:n =
      {
        \bool_gset_false:N \g__hit_bilingual_caption_auto_bool
        \bool_gset_true:N \g__hit_bilingual_caption_bool
      } ,
    bilingual-caption / false .code:n =
      {
        \bool_gset_false:N \g__hit_bilingual_caption_auto_bool
        \bool_gset_false:N \g__hit_bilingual_caption_bool
      } ,
    bilingual-caption / auto .code:n =
      { \bool_gset_true:N \g__hit_bilingual_caption_auto_bool } ,
    bilingual-caption .default:n = { true } ,
  }
%    \end{macrocode}
%
% 排双语题注：走 \pkg{ccaption} 的原形，英文图表名取 \option{const} 族里的
% \option{figure-prefix-en} 与 \option{table-prefix-en}，用户不必每次手写
% \texttt{Fig.} 加负细空格去抵 \cs{fnum@figure} 插的那个空格。
%
% 短标题给 \cs{bicaption} 的\emph{第二个必选参数}，不是可选参数。ccaption 的
% \cs{bicaption}\oarg{标签}\marg{短标题}\marg{题注}\marg{英文图表名}\marg{英文题注}
% 里，可选参数是 \cs{label}：
% \begin{latex}
% \newcommand{\bicaption}[5][\@empty]{%
%   \@if@contemptyarg{#2}{\caption{#3}}{\caption[#2]{#3}}
%   \ifx\@empty#1\else \label{#1} \fi
% \end{latex}
% 接错了的后果是同一个 \cs{caption}\oarg{短标题} 两条路径含义不同：单语时进插图索引，
% 双语时变成一个标签，索引里排的还是完整题注。
%
% \pkg{ccaption} 之所以提供“可选参数当标签”这条路，是因为 \cs{bicaption} 把整段
% 包在 \cs{begingroup}\ldots\cs{endgroup} 里，\cs{caption} 设的 \cs{@currentlabel}
% 是局部的，出了组就没了——写在后面的 \cs{label} 只能拿到章节号。所以排完之后
% 按计数器重建一次 \cs{@currentlabel}，让 \cs{label} 与普通 \hologo{LaTeX} 一样写在
% \cs{caption} 后面。
%    \begin{macrocode}
%    \end{macrocode}
%
% \emph{英文那一行按西文的行盒排。} 整行只有西文，Word 的行盒就按 Times 算
% （单倍 1.15 倍字号），跟着题注目标写的宋体算会高出一截：五号 1.25 倍下
% 两行的步进 15.2 变成 16.2。\pkg{ccaption} 在两条题注之间留了 \cs{midbicaption}
% 这个钩子，在那里立 \cs{l\_docxlike\_format\_latin\_pick\_bool}，与目录里整条西文的
% 条目同一个机关（\file{src/engine/hit-keylist.dtx} 的 \cs{hit@toc@latin@pick:n}）。
% 标志位是局部的，\cs{bicaption} 自己套着一层组，出组即还。
%    \begin{macrocode}
\cs_new_protected:Npn \hit@caption@latin@pick:
  {
    \bool_if_exist:NT \l_docxlike_format_latin_pick_bool
      { \bool_set_true:N \l_docxlike_format_latin_pick_bool }
  }
\cs_new_protected:Npn \hit@bilingual@caption #1#2#3
  {
    \midbicaption { \hit@caption@latin@pick: }
    \tl_if_novalue:nTF {#1}
      { \__hit_bicaption:w {} {#2} { \use:c { hit@term@caption@ \@captype @prefix@en } } {#3} }
      { \__hit_bicaption:w {#1} {#2} { \use:c { hit@term@caption@ \@captype @prefix@en } } {#3} }
    \protected@edef \@currentlabel
      { \use:c { p@ \@captype } \use:c { the \@captype } }
  }
\cs_new_protected:Npn \hit@single@caption #1#2
  {
    \tl_if_novalue:nTF {#1}
      { \hit@original@caption {#2} }
      { \hit@original@caption [#1] {#2} }
  }
\bool_new:N \l_hit_caption_star_bool
\cs_new_protected:Npn \hit@star@caption #1#2
  {
    \group_begin:
      \bool_set_true:N \l_hit_caption_star_bool
      \hit@original@caption * {#1}
      \bool_lazy_and:nnT
        { ! \tl_if_novalue_p:n {#2} }
        { \g__hit_bilingual_caption_bool }
        {
          \hit@caption@latin@pick:
          \hit@original@caption * {#2}
        }
    \group_end:
  }
%    \end{macrocode}
%
% \begin{macro}{\hit@float@sep:}
% \option{style}/\option{float-sep} 给了长度，三个间距一律照它，不给伸缩。
% 学位论文那边按网格行设的那一处（\file{src/flow/hit-final-float.dtx}）先看这个
% 键有没有值，有就让开。
%    \begin{macrocode}
\cs_new_protected:Npn \hit@float@sep:
  {
    \tl_if_empty:NF \g__hit_float_sep_tl
      {
        \skip_set:Nn \intextsep    { \g__hit_float_sep_tl }
        \skip_set:Nn \textfloatsep { \g__hit_float_sep_tl }
        \skip_set:Nn \floatsep     { \g__hit_float_sep_tl }
      }
  }
%    \end{macrocode}
% \end{macro}
%
% \cs{caption} 的接管。给了尾随的英文题、并且开关没关死，才走双语。
%
% 星号版交还原来的 \cs{caption*}，不编号、不进目录，续图续表的“图 1-1（续图）”
% 就靠它。签名里不收星号的话，\verb|\caption*{...}| 会把星号当成题注正文排出
% 一条“图 1-2 *”，后面的花括号组落成一段正文。星号版也认尾随的英文题，双语开着
% 就在下面再排一条不编号的。原来那个 \cs{caption*} 到了类自己的 \cs{@makecaption}
% 手里照样带着“图 1-1”的标签，所以这里立一个标志位，\cs{@makecaption} 见了就只排
% 题注正文（\file{hit-final-float.dtx}）。
%    \begin{macrocode}
\cs_new_protected:Npn \hit@setup@caption
  {
    \cs_set_eq:NN \__hit_bicaption:w \bicaption
    \AddToHook { begindocument/end }
      {
    \hit@float@sep:
    \bool_if:NT \g__hit_bilingual_caption_auto_bool
      {
        \bool_gset_eq:NN \g__hit_bilingual_caption_bool \g__hit_doctor_bool
        \bool_if_exist:NT \g__hit_postdoc_bool
          {
            \bool_if:NT \g__hit_postdoc_bool
              { \bool_gset_true:N \g__hit_bilingual_caption_bool }
          }
      }
    \cs_set_eq:NN \hit@original@caption \caption
    \RenewDocumentCommand \caption { s o +m o }
      {
        \IfBooleanTF {##1}
          { \hit@star@caption {##3} {##4} }
          {
            \IfNoValueTF {##4}
              { \hit@single@caption {##2} {##3} }
              {
                \bool_if:NTF \g__hit_bilingual_caption_bool
                  { \hit@bilingual@caption {##2} {##3} {##4} }
                  { \hit@single@caption {##2} {##3} }
              }
          }
      }
%    \end{macrocode}
%
% \cs{bicaption} 的兼容层。三种写法靠“后面还跟不跟一个花括号组”区分：跟着的是
% \pkg{ccaption} 的五参数原形，抓走那一组照原形排；不跟就是本模板之前的三、四参数
% 写法。探之前先跳空格，与 \cs{@ifnextchar} 的惯例一致——旧示例里那五个参数是连着
% 写的，但用户可能换了行。题注后面正常不会紧跟花括号组。
%    \begin{macrocode}
    \RenewDocumentCommand \bicaption { o +m +m +g }
      {
        \peek_meaning_ignore_spaces:NTF \c_group_begin_token
          { \__hit_bicaption_legacy:nnnn {##1} {##2} {##3} {##4} }
          {
            \msg_warning:nn { hithesis } { bicaption-deprecated }
            \IfNoValueTF {##4}
              { \hit@bilingual@caption {##1} {##2} {##3} }
              { \hit@bilingual@caption {##1} {##3} {##4} }
          }
      }
      }
  }
\cs_new_protected:Npn \__hit_bicaption_legacy:nnnn #1#2#3#4#5
  {
    \msg_warning:nn { hithesis } { bicaption-five-args }
    \tl_if_novalue:nTF {#1}
      { \__hit_bicaption:w {#2} {#3} {#4} {#5} }
      { \__hit_bicaption:w [#1] {#2} {#3} {#4} {#5} }
  }
%    \end{macrocode}
% \end{macro}
%
% \begin{macro}{\figurenote}
% \cs{figurenote}\marg{文字} 排图注。规范里图注是图的一部分，位置写死在图题之上：
% \begin{quote}
% 图中文字用宋体 5 号字\ldots 有图注或其它说明时应置于图题之上。
% \end{quote}
% 所以写在 \cs{caption} 前面。字号行距从 \option{styles} 表的 \option{caption-note}
% 取，与题注同一档。
%
% \emph{为什么值得给一个命令。} 让用户自己写 \cs{minipage} 加 \cs{vskip} 加
% \cs{wuhao}，同一份文档里两处图注就会两套数，改一处不动另一处。给了命令之后
% 格式只有一个出处，用户写内容就行。
%
% 打头那个 \cs{par} 是必要的：\cs{docxlike\_format\_par:nn} 里有 \cs{nointerlineskip}，
% 那条只能在竖排模式下用。图注常常紧跟在 \cs{end}\marg{minipage} 或
% \cs{includegraphics} 后面，那时还在横排模式，不先断段会报
% \texttt{You~can't~use~\string\prevdepth~in~horizontal~mode}。
%    \begin{macrocode}
\NewDocumentCommand \figurenote { m }
  { \par \docxlike_format_par:nn { caption-note } {#1} }
%    \end{macrocode}
% \end{macro}
%
% \emph{题号与题名之间的间隔。} 图题表题各读一个常量，常量为空退回兜底值：
% 新版面一律一个汉字宽（\cs{ccwd}，题注贴字符网格，一个字就是一格）；旧版面
% 照 v3.1e，本科一个字宽、其余 \cs{enskip}。图还是表看 \cs{@captype}；
% \cs{LT@makecaption} 那条路上它也是有的，取不到就按图算。
%
% \emph{为什么在共享层。} \file{src/engine/hit-table.dtx} 里给 \pkg{tabularray} 长表
% 注册 \option{caption-en} 外参那一段是 \cs{cs\_if\_exist:NT} \cs{hit@caption@separator}
% 守着的，报告里取不到这个宏，那个外参就没注册，照抄一张终稿的长表会报
% “Unknown~outer~key~name~'caption-en'”。五条本身没有一处与终稿绑定。
%    \begin{macrocode}
\cs_new:Npn \hit@caption@separator@fallback
  {
    \dim_compare:nNnTF { \g_docxlike_page_lead_dim } > { \c_zero_dim }
      { \hskip \ccwd }
      { \hit@if@option@TF { bachelor } { \hskip \ccwd } { \enskip } }
  }
\cs_new:Npn \hit@caption@number@lang:
  { \bool_if:NTF \g__hit_en_bool { en } { zh } }
\cs_new:Npn \hit@caption@type:
  { \cs_if_exist:NTF \@captype { \@captype } { figure } }
%    \end{macrocode}
%
% \cs{hit@caption@font:} 按图还是表挑目标，从 \option{styles} 表取字号、行距、
% 贴不贴网格。表里没取值就退回 \cs{wuhao}——旧版面那一份表是空的，行距跟着
% \cs{wuhao} 自带的 1.3 倍走，与 v3.1e 一样。
%    \begin{macrocode}
\cs_new_protected:Npn \hit@caption@font:
  {
    \docxlike_format_apply_font_else:en { caption- \hit@caption@type: } { \wuhao }
  }
\cs_new:Npn \hit@caption@separator
  {
    \str_if_eq:eeTF { \hit@caption@type: } { table }
      {
        \hit@after@number:nnn { table@caption } { \hit@caption@number@lang: }
          { \hit@caption@separator@fallback }
      }
      {
        \hit@after@number:nnn { figure@caption } { \hit@caption@number@lang: }
          { \hit@caption@separator@fallback }
      }
  }
%</shared-caption>
%    \end{macrocode}
%
% \Finale
